\documentclass[11pt]{article}
\usepackage{kctdocs}
\begin{document}
\kcttitle{}{Transcription Style Guide}{%
Status & v0.2: decisions confirmed; examples to be extended in Phase 0 (P0.8) \\
Version & 1.2 \textperiodcentered\ 2026-09-24 22:26 \textperiodcentered\ \hyperref[changelog]{change log} \\
Applies to & Source transcripts (DOCX), reviewer edits in Corpus Studio, and training targets \\
}
{\small\setlength{\parskip}{0pt}\setcounter{tocdepth}{1}\tableofcontents}
\bigskip
\hypertarget{two-layers-of-text}{%
\section{Two layers of text}\label{two-layers-of-text}}

\begin{longtable}[]{@{}>{\raggedright\arraybackslash}p{(\linewidth - 8\tabcolsep) * \real{0.105}}>{\raggedright\arraybackslash}p{(\linewidth - 8\tabcolsep) * \real{0.178}}>{\raggedright\arraybackslash}p{(\linewidth - 8\tabcolsep) * \real{0.396}}>{\raggedright\arraybackslash}p{(\linewidth - 8\tabcolsep) * \real{0.321}}@{}}
\toprule\noalign{}
Layer & Where & Numbers & Purpose \\
\midrule\noalign{}
\endhead
\bottomrule\noalign{}
\endlastfoot
\textbf{Source text} & The DOCX transcripts & Digits (\texttt{6}, \texttt{2019}, \texttt{Level 5}) & Easy for people to write and read. Never modified by the pipeline \\
\textbf{Spoken form} & Segment revisions → training targets & Words in the language actually spoken (\texttt{sita} or \texttt{six}, \texttt{twenty nineteen}, \texttt{Level Five}) & What the model hears and learns \\
\textbf{Display form} & Transcribe output (optional) & Digits again, via inverse text normalisation & Readable subtitles and documents \\
\end{longtable}

The pipeline converts source text to spoken form automatically (§4). You keep writing digits in the documents.

\hypertarget{language-marking-in-docx}{%
\section{Language marking in DOCX}\label{language-marking-in-docx}}

\begin{itemize}
\tightlist
\item
  \textbf{Italic = Kiswahili}; roman (upright) = English. This applies to every document.
\item
  \textbf{Hybrid words:} write them as \textbf{one word, no hyphen}. Italicise only the Kiswahili part if you can: \emph{nime}download, \emph{wame}organize, \emph{tuta}text\emph{iana}. The parser tags a word that is part italic and part roman as \texttt{mixed}. If splitting the formatting is too slow, italicising the whole hybrid word is acceptable; the reviewer can retag it.
\item
  Sheng words written with Kiswahili grammar are italic. English loanwords already absorbed into Kiswahili spelling (\emph{kompyuta}, \emph{basi}) are italic.
\item
  Words that are neither English nor Kiswahili (e.g. Dholuo): write them as heard and wrap them in \texttt{[luo:\allowbreak{} …]}, e.g. \texttt{[luo:\allowbreak{} ber ahinya]}.
\end{itemize}

\hypertarget{verbatim-rules}{%
\section{Verbatim rules}\label{verbatim-rules}}
\kctchanged{1.2}

\begin{itemize}
\tightlist
\item
  Transcribe what was said, including repetitions that carry meaning and every code-switch.
\item
  Drop pure fillers (\emph{eh}, \emph{umm}, \emph{aah}) unless they carry meaning (\emph{eeh} = ``yes'').
\item
  Unintelligible speech: \texttt{[Unclear 13:10-13:15]}, with the start and end time (\texttt{h:mm:ss} after the first hour: \texttt{[Unclear 1:05:21-1:05:24]}); add a note if useful: \texttt{[Unclear - Luo, 37:39-38:45]}. Do not guess. These segments are excluded from training until resolved.
\item
  Long stretches without transcribable speech: \texttt{[No speech - activity, 17:26-18:14]}.
\item
  \textbf{Always include the times} when a marked stretch is longer than about 2 seconds. The aligner uses them as fixed points, so words cannot drift into the stretch, and the report uses them to measure alignment accuracy. Use square brackets only: the tools also accept \texttt{(Unclear: 1:16:07-1:16:08]]}, but mixed brackets are easy to get wrong.
\item
  Participant references (\texttt{P7}, \texttt{P12}) stand in for a participant's name for anonymity. They are \emph{not} read as ``pee seven''. Chunks containing one are kept out of training, because the audio holds a name the text does not.
\item
  False starts and cut-off words: \texttt{tuta--tutaanza}, \texttt{tuna--}. The double hyphen marks the break and the parts are treated as separate words.
\item
  Crosstalk: \texttt{[overlap]} at the start of the overlapping part.
\item
  Passages you paraphrased or skipped for meaning: \texttt{[nv]} \ldots{} \texttt{[/\allowbreak{}nv]} (non-verbatim). These segments are excluded from training until corrected in review.
\end{itemize}

\hypertarget{numbers}{%
\section{Numbers}\label{numbers}}

Source documents always use \textbf{digits}. The spoken form uses \textbf{words in the language actually spoken}.

\hypertarget{default-language-rules-applied-automatically}{%
\subsection{Default language rules (applied automatically)}\label{default-language-rules-applied-automatically}}

\begin{longtable}[]{@{}>{\raggedright\arraybackslash}p{(\linewidth - 6\tabcolsep) * \real{0.374}}>{\raggedright\arraybackslash}p{(\linewidth - 6\tabcolsep) * \real{0.311}}>{\raggedright\arraybackslash}p{(\linewidth - 6\tabcolsep) * \real{0.316}}@{}}
\toprule\noalign{}
Case & Rule & Source → spoken \\
\midrule\noalign{}
\endhead
\bottomrule\noalign{}
\endlastfoot
Number inside an \emph{italic} (Kiswahili) span & Kiswahili words & \emph{watu 6} → watu sita \\
Number inside a roman (English) span & English words & 6 people → six people \\
Year (4 digits, 1900--2099) & English words, year style & 2019 → twenty nineteen; 2005 → two thousand and five \\
Label + number (Level, Grade, Form, Class, Chapter, Room, Route, Phase\ldots) & English words, title case & Level 5 → Level Five; Form 4 → Form Four \\
Money & Follow the span language; the currency word stays as spoken & \emph{shilingi 500} → shilingi mia tano; 500 bob → five hundred bob \\
Phone and ID numbers & Digit by digit, span language & 0712 → zero seven one two \\
Time & Follow the span language & \emph{saa 3} → saa tatu; 3 pm → three pm \\
\end{longtable}

\hypertarget{why-a-rule-alone-isnt-enough-and-how-it-gets-resolved}{%
\subsection{Why a rule alone isn\textquotesingle t enough, and how it gets resolved}\label{why-a-rule-alone-isnt-enough-and-how-it-gets-resolved}}

A bare \texttt{6} can\textquotesingle t tell you whether the speaker said \emph{sita} or \emph{six}. Kenyans often say numbers, and especially years and money, in English inside Kiswahili sentences. So the pipeline doesn\textquotesingle t trust the rule blindly:

\begin{enumerate}
\def\labelenumi{\arabic{enumi}.}
\tightlist
\item
  It generates both candidates (\texttt{sita} and \texttt{six}; \texttt{twenty nineteen} and \texttt{elfu mbili kumi na tisa}).
\item
  The forced aligner scores each candidate against the audio, and the higher-scoring one wins.
\item
  If the two scores are too close, the segment gets a \texttt{number\_\allowbreak{}check} flag and the reviewer hears it and picks.
\end{enumerate}

In practice the audio decides, and the §4.1 table is only the tie-breaker.

\hypertarget{kiswahili-number-words-reference}{%
\subsection{Kiswahili number words (reference)}\label{kiswahili-number-words-reference}}

1 moja · 2 mbili · 3 tatu · 4 nne · 5 tano · 6 sita · 7 saba · 8 nane · 9 tisa · 10 kumi · 11 kumi na moja · 20 ishirini · 30 thelathini · 40 arobaini · 50 hamsini · 60 sitini · 70 sabini · 80 themanini · 90 tisini · 100 mia moja · 250 mia mbili hamsini · 1,000 elfu moja · 2019 elfu mbili kumi na tisa · 1,000,000 milioni moja

Noun-class agreement (\emph{watu watatu}, \emph{vitu vitatu}) is written \textbf{as spoken} by the reviewer. The automatic converter emits the bare counting form (\emph{tatu}), and the aligner or reviewer corrects it to the agreeing form.

\hypertarget{other-spelling-conventions}{%
\section{Other spelling conventions}\label{other-spelling-conventions}}

\begin{longtable}[]{@{}>{\raggedright\arraybackslash}p{(\linewidth - 4\tabcolsep) * \real{0.138}}>{\raggedright\arraybackslash}p{(\linewidth - 4\tabcolsep) * \real{0.862}}@{}}
\toprule\noalign{}
Item & Convention \\
\midrule\noalign{}
\endhead
\bottomrule\noalign{}
\endlastfoot
Casing & Sentence case; proper nouns capitalised in both languages \\
Punctuation & Normal punctuation in transcripts (Whisper learns it); ignored when scoring \\
Contractions & As spoken: \emph{don\textquotesingle t}, \emph{I\textquotesingle m}; Kiswahili elisions as spoken (\emph{nshaenda}) \\
Okay & \emph{okay} (not \emph{ok}, \emph{OK}, \emph{sawa} unless \emph{sawa} was said) \\
Acronyms & Uppercase, no dots: \emph{KCSE}, \emph{NHIF}; letter-by-letter acronyms stay uppercase \\
\end{longtable}

\kctchangelog
\begin{kctversions}
1.0 & 2026-09-24 00:36 & \texttt{053e09d} & First LaTeX version (v0.1 of the guide). \\
1.2 & 2026-09-24 22:26 & this commit & Marker format with timestamps, participant references, false starts (from the NRCCW\_KSM09 pilot). \\
\end{kctversions}

\kctchange{1.2}{§3 Verbatim rules}{verbatim-rules}
 {\kctsrc{8a3be55a33b1a843a27d87795a74872523d940bb}{STYLE_GUIDE.tex}{43}{57}{v1.1 source, lines 43--57} \textperiodcentered\ \href{archive/v1.1/STYLE_GUIDE.pdf}{v1.1 PDF}}
 {``Unintelligible speech: \texttt{[unclear]}. Do not guess.''; nothing on participant references, no-speech stretches or false starts.}
 {Markers carry start and end times (\texttt{[Unclear 13:10-13:15]}, \texttt{[No speech - activity, 17:26-18:14]}), which the aligner uses as fixed points; square brackets only. \texttt{P7}-style participant references are anonymised names, not spoken letters, and keep their chunk out of training. False starts use \texttt{--}.}
\end{document}
